% SPDX-License-Identifier: GPL-3.0-or-later
% Copyright (C) 2026 Alessandro Greco
%% 简体中文版 EPUB 替代文本。
\altfig{Chaum 式混合节点在左侧接收四条等长加密消息，标记为 A、B、C、D。等待凑齐一批、去除一层加密并打乱顺序后，右侧输出四条统一外观的消息，标记为 1 至 4；图中不显示任何输入与输出的对应关系。}
% Chaum 的 mix：批处理、去除一层加密并打乱顺序。
\begin{tikzpicture}[font=\sffamily\scriptsize,
  msg/.style={draw,minimum width=1.5cm,minimum height=0.5cm,rounded corners=1pt}]
  \node[draw=torviola,fill=torviola!8,rounded corners=4pt,minimum width=2.4cm,minimum height=3.4cm] (mix) at (5.2,0) {};
  \node[font=\sffamily\small,text=torviola] at (5.2,1.35) {\textbf{mix}};
  \node[align=center,text=torviola] at (5.2,0.4) {等待成批\\去除一层\\打乱顺序};
  \foreach \i/\y/\c in {A/1.2/torverde!30,B/0.4/torviola!25,C/-0.4/avviso!25,D/-1.2/torgrigio!25} {
    \node[msg,fill=\c] (in\i) at (1.6,\y) {\i};
    \node[font=\sffamily\tiny,left=1pt of in\i] {发送者 \i};
  }
  \foreach \i/\y in {1/1.2,2/0.4,3/-0.4,4/-1.2} {
    \node[msg,fill=torgrigio!20] (out\i) at (8.8,\y) {\i};
  }
  \foreach \i in {A,B,C,D} \draw[->,torgrigio] (in\i) -- (mix.west |- in\i);
  \foreach \i in {1,2,3,4} \draw[->,torgrigio] (mix.east |- out\i) -- (out\i);
  \node[align=center,font=\sffamily\tiny] at (1.6,-1.9) {长度相同，\\为 mix 加密};
  \node[align=center,font=\sffamily\tiny] at (8.8,-1.9) {外观已改变，\\顺序已打乱};
  \draw[decorate,decoration={brace,amplitude=4pt}] (0.6,1.55) -- node[above=4pt] {输入批次} (2.5,1.55);
  \draw[decorate,decoration={brace,amplitude=4pt}] (7.9,1.55) -- node[above=4pt] {输出批次} (9.7,1.55);
\end{tikzpicture}
